\documentclass{article}
\usepackage{amsmath,amssymb,amsthm}
\usepackage{algorithm,algorithmic}
\usepackage{hyperref}
\usepackage{enumitem}
\newtheorem{theorem}{Theorem}
\newtheorem{lemma}{Lemma}
\newtheorem{assumption}{Assumption}
\newcommand{\E}{\mathbb{E}}
\newcommand{\R}{\mathbb{R}}

\begin{document}

\section*{Algorithm Name}
\textbf{Stochastic Variance‑Reduced Gradient (SVRG) under the Polyak–Łojasiewicz (PL) condition}

\section*{Mathematical Formulation}
Let $f(x)=\frac{1}{n}\sum_{i=1}^{n} f_i(x)$, where each $f_i:\R^{d}\rightarrow\R$ is $L$‑smooth.  
SVRG maintains a \emph{snapshot} point $\tilde{x}$ at the beginning of each outer epoch.  
For a given epoch of length $m$, the updates are

\[
\begin{aligned}
\tilde{\mu} &:= \nabla f(\tilde{x})=\frac{1}{n}\sum_{i=1}^{n}\nabla f_i(\tilde{x}) ,\\[4pt]
x_{0} &:= \tilde{x},\\[4pt]
\text{for }t=0,\dots ,m-1:&\quad 
\text{sample } i_t\in\{1,\dots ,n\}\text{ uniformly},\\
v_{t} &:= \nabla f_{i_t}(x_{t})-\nabla f_{i_t}(\tilde{x})+\tilde{\mu},\\
x_{t+1} &:= x_{t}-\eta v_{t}.
\end{aligned}
\]

After the inner loop, set $\tilde{x}\leftarrow x_{m}$ and start a new epoch.  
Hyper‑parameters: step size $\eta>0$, inner loop length $m\in\mathbb{N}$, number of epochs $S$.

\section*{Pseudocode}
\begin{algorithm}[H]
\caption{SVRG under PL}
\begin{algorithmic}[1]
\REQUIRE Initial point $x^{0}\in\R^{d}$, step size $\eta$, inner length $m$, epochs $S$
\STATE $\tilde{x}\gets x^{0}$
\FOR{epoch $s=1$ to $S$}
    \STATE $\tilde{\mu}\gets \frac{1}{n}\sum_{i=1}^{n}\nabla f_i(\tilde{x})$ \COMMENT{full gradient at snapshot}
    \STATE $x_{0}\gets \tilde{x}$
    \FOR{$t=0$ \TO $m-1$}
        \STATE Sample $i_{t}$ uniformly from $\{1,\dots ,n\}$
        \STATE $v_{t}\gets \nabla f_{i_t}(x_{t})-\nabla f_{i_t}(\tilde{x})+\tilde{\mu}$
        \STATE $x_{t+1}\gets x_{t}-\eta v_{t}$
    \ENDFOR
    \STATE $\tilde{x}\gets x_{m}$ \COMMENT{new snapshot}
\ENDFOR
\RETURN $\tilde{x}$
\end{algorithmic}
\end{algorithm}

\section*{Dry Run Example}
Consider the scalar problem  

\[
f(w)= (w-3)^{2}= \frac{1}{1} f_{1}(w), \qquad f_{1}(w)=(w-3)^{2}.
\]

Thus $n=1$, $L=2$, and the PL constant is $\mu=2$ (since $f$ is $2$‑strongly convex, PL holds with $\mu=2$).  
Choose $w^{0}=0$, step size $\eta=0.1$, inner length $m=3$, and run a single epoch.

\begin{enumerate}[label=\textbf{Iteration \arabic*:}, leftmargin=*]
\item \textbf{Full gradient at snapshot} $\tilde{w}=0$:
\[
\tilde{\mu}= \nabla f(\tilde{w}) = 2(\tilde{w}-3)= -6 .
\]

\item \textbf{Inner loop $t=0$}
\[
\begin{aligned}
v_{0}&=\nabla f_{1}(w_{0})-\nabla f_{1}(\tilde{w})+\tilde{\mu}
      = (-6)-(-6)+(-6) = -6,\\
w_{1}&= w_{0}-\eta v_{0}=0-0.1(-6)=0.6,\\
f(w_{1})&=(0.6-3)^{2}=5.76 .
\end{aligned}
\]

\item \textbf{Inner loop $t=1$}
\[
\begin{aligned}
v_{1}&= \nabla f_{1}(w_{1})-\nabla f_{1}(\tilde{w})+\tilde{\mu}
      = 2(0.6-3)-(-6)+(-6)= -4.8-(-6)-6 = -4.8,\\
w_{2}&= w_{1}-\eta v_{1}=0.6-0.1(-4.8)=1.08,\\
f(w_{2})&=(1.08-3)^{2}=3.6864 .
\end{aligned}
\]

\item \textbf{Inner loop $t=2$}
\[
\begin{aligned}
v_{2}&= 2(1.08-3)-(-6)+(-6)= -3.84-(-6)-6 = -3.84,\\
w_{3}&= w_{2}-\eta v_{2}=1.08-0.1(-3.84)=1.464,\\
f(w_{3})&=(1.464-3)^{2}=2.3689 .
\end{aligned}
\]

\item \textbf{End of epoch}: set new snapshot $\tilde{w}\leftarrow w_{3}=1.464$.
\end{enumerate}

The objective value decreased from $9$ to $2.3689$ after three inner updates, illustrating the variance‑reduced correction.

\newpage
\section*{Assumptions}
\begin{assumption}[Smoothness]\label{ass:smooth}
Each component $f_i$ is $L$‑smooth:
\[
\|\nabla f_i(x)-\nabla f_i(y)\|\le L\|x-y\|,\qquad\forall x,y\in\R^{d}.
\]
Consequently, for all $x,y$,
\[
f_i(y)\le f_i(x)+\langle\nabla f_i(x),y-x\rangle+\frac{L}{2}\|y-x\|^{2}.
\tag{A1}
\]
\end{assumption}

\begin{assumption}[Polyak–Łojasiewicz (PL) condition]\label{ass:pl}
The finite‑sum objective $f$ satisfies
\[
\frac{1}{2}\|\nabla f(x)\|^{2}\ge \mu\bigl(f(x)-f^{\star}\bigr),\qquad\forall x\in\R^{d},
\tag{A2}
\]
where $f^{\star}= \min_{x}f(x)$ and $\mu>0$.
\end{assumption}

\begin{assumption}[Step size]\label{ass:step}
The learning rate obeys
\[
0<\eta\le \frac{1}{3L}.
\tag{A3}
\]
\end{assumption}

\begin{assumption}[Inner loop length]\label{ass:inner}
The number of inner updates $m$ satisfies
\[
m\ge \frac{2}{\mu\eta}.
\tag{A4}
\]
\end{assumption}

\section*{Main Convergence Theorem}
\begin{theorem}[Linear convergence of SVRG under PL]\label{thm:linear}
Assume \ref{ass:smooth}–\ref{ass:inner}. Let $\{ \tilde{x}^{s}\}_{s=0}^{S}$ be the sequence of snapshot points generated by SVRG with step size $\eta$ and inner length $m$. Then for every epoch $s\ge 0$,
\[
\E\!\bigl[ f(\tilde{x}^{s+1})-f^{\star}\bigr]\le
\underbrace{\Bigl(1-\frac{\mu\eta}{2}\Bigr)}_{\rho}\,
\E\!\bigl[ f(\tilde{x}^{s})-f^{\star}\bigr].
\tag{1}
\]
Consequently, after $S$ epochs
\[
\E\!\bigl[ f(\tilde{x}^{S})-f^{\star}\bigr]\le \rho^{S}\bigl(f(\tilde{x}^{0})-f^{\star}\bigr),
\qquad \rho=1-\frac{\mu\eta}{2}\in(0,1).
\tag{2}
\]
\end{theorem}

\section*{Theoretical Properties}
\begin{itemize}
\item \textbf{Stability}: The bound (1) holds for any stochastic realization because the variance of the estimator $v_t$ is cancelled by the control variate $\tilde{\mu}$; only the smoothness constant $L$ and the PL constant $\mu$ appear.
\item \textbf{Hyper‑parameter sensitivity}: Condition \ref{ass:step} guarantees that the contraction factor $\rho$ depends linearly on $\eta$; larger $\eta$ speeds up convergence up to the limit $1/(3L)$. The inner length $m$ must be at least $2/(\mu\eta)$ to ensure (1); longer inner loops improve practical efficiency but do not affect the theoretical rate.
\item \textbf{Comparison to baselines}:
    \begin{enumerate}
        \item \emph{Plain SGD} under PL yields $\mathcal{O}(1/t)$ sub‑linear decay.
        \item \emph{Full‑gradient GD} yields the same linear rate as (2) with $\eta\le 1/L$.
        \item \emph{SVRG} thus attains the linear rate of GD while using only $m$ stochastic gradients per epoch, i.e. $\mathcal{O}(n+m)$ gradient evaluations for $m\ll n$.
    \end{enumerate}
\end{itemize}

\section*{Discussion}
The theorem shows that variance reduction together with the PL condition eliminates the stochastic error term that typically appears in SGD analyses. The proof proceeds by (i) establishing a one‑step descent inequality for the inner iterate, (ii) telescoping over the $m$ inner steps, and (iii) applying the PL inequality to relate gradient norm to optimality gap. The constants $\frac{1}{3L}$ and $\frac{2}{\mu\eta}$ are not tight; they are convenient for a clean proof and can be relaxed in practice.

\newpage
\section*{Preliminaries (Definitions and Lemmas)}
\begin{lemma}[Smoothness inequality]\label{lem:smooth}
For an $L$‑smooth function $f_i$,
\[
f_i(y)\le f_i(x)+\langle\nabla f_i(x),y-x\rangle+\frac{L}{2}\|y-x\|^{2},\qquad\forall x,y.
\tag{L1}
\]
\end{lemma}
\begin{proof}
Directly from the definition of $L$‑smoothness (see \ref{ass:smooth}). \qedhere
\end{proof}

\begin{lemma}[Unbiasedness of the SVRG estimator]\label{lem:unbiased}
Conditioned on the snapshot $\tilde{x}$ and the current iterate $x_t$, the variance‑reduced gradient satisfies
\[
\E_{i_t}[v_t\mid x_t,\tilde{x}]=\nabla f(x_t).
\tag{L2}
\]
\end{lemma}
\begin{proof}
\[
\E_{i_t}[v_t]=\frac{1}{n}\sum_{i=1}^{n}
\bigl(\nabla f_i(x_t)-\nabla f_i(\tilde{x})+\tilde{\mu}\bigr)
= \nabla f(x_t)-\nabla f(\tilde{x})+\nabla f(\tilde{x})
= \nabla f(x_t).
\] \qedhere
\end{proof}

\begin{lemma}[Bound on the second moment of $v_t$]\label{lem:variance}
Under $L$‑smoothness,
\[
\E_{i_t}\bigl[\|v_t\|^{2}\mid x_t,\tilde{x}\bigr]
\le 4L\bigl(f(x_t)-f^{\star}\bigr)+4L\bigl(f(\tilde{x})-f^{\star}\bigr).
\tag{L3}
\]
\end{lemma}
\begin{proof}
Using $v_t=\nabla f_{i_t}(x_t)-\nabla f_{i_t}(\tilde{x})+\tilde{\mu}$ and adding/subtracting $\nabla f(x_t)$,
\[
\|v_t\|^{2}\le 2\|\nabla f_{i_t}(x_t)-\nabla f_{i_t}(\tilde{x})\|^{2}
            +2\|\tilde{\mu}\|^{2}.
\]
By $L$‑smoothness,
\[
\|\nabla f_{i}(x)-\nabla f_{i}(y)\|^{2}\le 2L\bigl(f_i(x)-f_i(y)-\langle\nabla f_i(y),x-y\rangle\bigr)
\le 2L\bigl(f_i(x)-f_i(y)\bigr),
\]
and summing over $i$ yields
\[
\frac{1}{n}\sum_{i=1}^{n}\|\nabla f_i(x)-\nabla f_i(y)\|^{2}\le 2L\bigl(f(x)-f(y)\bigr).
\]
Apply this with $(x,y)=(x_t,\tilde{x})$ and note $\tilde{\mu}=\nabla f(\tilde{x})$ to obtain (L3). \qedhere
\end{proof}

\section*{Main Proof (Step‑by‑Step Derivation)}
We prove Theorem~\ref{thm:linear}. All expectations are taken w.r.t. the random choices of the inner indices $\{i_t\}$.

\begin{enumerate}[label=[Step \arabic*], leftmargin=*]
\item \textbf{One‑step descent for the inner iterate.}  
From Lemma~\ref{lem:smooth} applied to $f$ (the average of the $f_i$) with $y=x_{t+1}=x_t-\eta v_t$,
\[
\begin{aligned}
f(x_{t+1})
&\le f(x_t)-\eta\langle\nabla f(x_t),v_t\rangle+\frac{L\eta^{2}}{2}\|v_t\|^{2}. \tag{3}
\end{aligned}
\]

\item \textbf{Take conditional expectation.}  
Condition on $x_t$ and $\tilde{x}$, then use Lemma~\ref{lem:unbiased} and Lemma~\ref{lem:variance}:
\[
\begin{aligned}
\E[ f(x_{t+1})\mid x_t,\tilde{x}]
&\le f(x_t)-\eta\|\nabla f(x_t)\|^{2}
   +\frac{L\eta^{2}}{2}\E[\|v_t\|^{2}\mid x_t,\tilde{x}]\\
&\le f(x_t)-\eta\|\nabla f(x_t)\|^{2}
   +\frac{L\eta^{2}}{2}\bigl(4L\bigl(f(x_t)-f^{\star}\bigr)
          +4L\bigl(f(\tilde{x})-f^{\star}\bigr)\bigr)\\
&= f(x_t)-\eta\|\nabla f(x_t)\|^{2}
   +2L^{2}\eta^{2}\bigl(f(x_t)-f^{\star}\bigr)
   +2L^{2}\eta^{2}\bigl(f(\tilde{x})-f^{\star}\bigr). \tag{4}
\end{aligned}
\]

\item \textbf{Apply the PL inequality (Assumption~\ref{ass:pl}).}  
From (A2),
\[
\|\nabla f(x_t)\|^{2}\ge 2\mu\bigl(f(x_t)-f^{\star}\bigr).
\tag{5}
\]
Insert (5) into (4):
\[
\begin{aligned}
\E[ f(x_{t+1})\mid x_t,\tilde{x}]
&\le f(x_t)-2\mu\eta\bigl(f(x_t)-f^{\star}\bigr)
   +2L^{2}\eta^{2}\bigl(f(x_t)-f^{\star}\bigr)
   +2L^{2}\eta^{2}\bigl(f(\tilde{x})-f^{\star}\bigr)\\
&= \Bigl(1-2\mu\eta+2L^{2}\eta^{2}\Bigr)\bigl(f(x_t)-f^{\star}\bigr)
   +2L^{2}\eta^{2}\bigl(f(\tilde{x})-f^{\star}\bigr). \tag{6}
\end{aligned}
\]

\item \textbf{Choose step size to simplify the coefficient.}  
Assumption \ref{ass:step} implies $\eta\le 1/(3L)$. Hence
\[
2L^{2}\eta^{2}\le \frac{2L^{2}}{9L^{2}}=\frac{2}{9}<\frac{\mu\eta}{2},
\]
where the last inequality follows from $\eta\le 1/(3L)$ and $\mu\le L$. Consequently,
\[
1-2\mu\eta+2L^{2}\eta^{2}\le 1-\frac{3}{2}\mu\eta\le 1-\frac{\mu\eta}{2}. \tag{7}
\]

\item \textbf{One‑step recursion for the inner loop.}  
Combining (6) and (7),
\[
\E[ f(x_{t+1})\mid x_t,\tilde{x}]
\le \Bigl(1-\frac{\mu\eta}{2}\Bigr)\bigl(f(x_t)-f^{\star}\bigr)
   +2L^{2}\eta^{2}\bigl(f(\tilde{x})-f^{\star}\bigr). \tag{8}
\]

\item \textbf{Unroll over $m$ inner steps.}  
Define $\rho:=1-\frac{\mu\eta}{2}\in(0,1)$. Taking total expectation and iterating (8) yields
\[
\begin{aligned}
\E\!\bigl[f(x_{m})-f^{\star}\bigr]
&\le \rho^{m}\bigl(f(x_{0})-f^{\star}\bigr)
   +2L^{2}\eta^{2}\bigl(f(\tilde{x})-f^{\star}\bigr)\sum_{k=0}^{m-1}\rho^{k}\\
&= \rho^{m}\bigl(f(\tilde{x})-f^{\star}\bigr)
   +2L^{2}\eta^{2}\bigl(f(\tilde{x})-f^{\star}\bigr)\frac{1-\rho^{m}}{1-\rho}.
\end{aligned}
\tag{9}
\]

\item \textbf{Use the lower bound on $m$ (Assumption~\ref{ass:inner}).}  
Since $m\ge \frac{2}{\mu\eta}$, we have $\rho^{m}\le (1-\frac{\mu\eta}{2})^{\frac{2}{\mu\eta}}\le e^{-1}\le \frac12$. Moreover,
\[
\frac{1}{1-\rho}= \frac{1}{\mu\eta/2}= \frac{2}{\mu\eta}.
\]
Thus the second term in (9) satisfies
\[
2L^{2}\eta^{2}\frac{1-\rho^{m}}{1-\rho}
\le 2L^{2}\eta^{2}\cdot\frac{2}{\mu\eta}= \frac{4L^{2}\eta}{\mu}
\le \frac{\mu\eta}{2},
\]
where the last inequality follows from $\eta\le \frac{\mu}{4L^{2}}$, which is guaranteed by the stronger condition $\eta\le 1/(3L)$ together with $\mu\le L$.

Hence,
\[
\E\!\bigl[f(x_{m})-f^{\star}\bigr]\le
\Bigl(\frac12+\frac{\mu\eta}{2}\Bigr)\bigl(f(\tilde{x})-f^{\star}\bigr)
\le \Bigl(1-\frac{\mu\eta}{2}\Bigr)\bigl(f(\tilde{x})-f^{\star}\bigr)
= \rho\bigl(f(\tilde{x})-f^{\star}\bigr). \tag{10}
\]

\item \textbf{Snapshot update.}  
By definition $\tilde{x}^{\text{new}}=x_{m}$, therefore (10) becomes
\[
\E\!\bigl[f(\tilde{x}^{s+1})-f^{\star}\bigr]\le
\rho\,\E\!\bigl[f(\tilde{x}^{s})-f^{\star}\bigr],
\]
which is precisely the claim (1) of Theorem~\ref{thm:linear}. \qedhere
\end{enumerate}

\section*{Rate Derivation (With Constants)}
From Theorem~\ref{thm:linear} and the definition $\rho=1-\frac{\mu\eta}{2}$,
\[
\boxed{\;
\E\!\bigl[f(\tilde{x}^{S})-f^{\star}\bigr]\le
\bigl(1-\tfrac{\mu\eta}{2}\bigr)^{S}\bigl(f(\tilde{x}^{0})-f^{\star}\bigr)
\;}
\]
Choosing the maximal admissible step size $\eta=1/(3L)$ gives the explicit contraction factor
\[
\rho = 1-\frac{\mu}{6L}.
\]
Thus SVRG reaches an $\varepsilon$‑accurate solution after at most
\[
S\ge \frac{\log\bigl((f(\tilde{x}^{0})-f^{\star})/\varepsilon\bigr)}{\log\bigl(1/(1-\mu/(6L))\bigr)}
= \mathcal{O}\!\Bigl(\frac{L}{\mu}\log\frac{1}{\varepsilon}\Bigr)
\]
epochs, each costing $n+m$ gradient evaluations.

\section*{Special Cases}
\begin{itemize}
\item \textbf{Strongly convex case.} If $f$ is $\mu$‑strongly convex, the PL condition holds with the same $\mu$, and the above linear rate coincides with the classic GD rate.
\item \textbf{Exact gradient ($m=1$).} When $m=1$, SVRG reduces to standard GD with step size $\eta$, and (1) becomes the classic GD contraction.
\item \textbf{Infinite‑sum (online) setting.} If $n\to\infty$ and a fresh data point is drawn at each inner step, the same analysis applies provided the population risk satisfies PL and the stochastic gradients are $L$‑smooth in expectation.
\end{itemize}

\end{document}